\chapter{任务六：确定性综合场景与可重复基准}
\label{chap:task6}

\section{约束}

综合任务既要连接至少三个模块，又要在真实 QEMU 中完成，还要提供性能对比。若把模型
随机性混入基准，重复运行的路径和工具数量可能不同；若只播放固定日志，又无法证明
机制真的参与。我们因此把“可重复基准”和“自由交互产品”分成两条同等重要但职责不同
的路径。

\code{contest-demo} 使用确定性 Guest policy，锁定业务工作量和 AB/BA 次序；Nexus
使用模型或 replay 动态委派，展示产品交互。前者回答“同负载下机制和路径如何表现”，
后者回答“真实用户目标如何被内核治理”。

\section{确定性科研恢复工作流}

\code{labdemo_ucore} 创建 Orchestrator、Sentinel、Investigator 和 Recovery，建立
workflow identity、MESSAGE route、Context、metadata 和 capability。场景围绕一个
科研任务运行：查询项目与 run 状态、发现失败 stage、读取依赖和摘要、发送跨 Agent
消息、尝试受限动作、形成恢复选择并复核结果。场景登记
\code{prepare/align/analyze/report/archive} 的状态和依赖，完成 align 恢复与 report 更新；
analyze 和 archive 保持 pending。

固定 policy 不是伪装成 LLM。代码显式安排阶段和参数，工具结果仍由内核真实执行并
检查；任何失败都会使场景失败。它适合验证身份、Context、文件索引、IPC wait、权限
拒绝和 recovery action 是否贯通，也适合在同一 AgentOS Guest 中对 traversal/compat 与
indexed/native 做业务合同一致的配对测量。

\code{labdemo_ucore} 输出的 \code quiescence 只用于稳定性能计数。真实
320-byte workflow fence receipt 由双目标路径中的 \code{rp_agentos_orch} 生成，并由
\code{make dual-platform-run} 复验。

\section{Native/Compat 与 AB/BA}

综合场景保留 traversal/compat 与 indexed/native 两条功能等价路径。每个独立 QEMU boot
内按 A/B 或 B/A 顺序运行，记录共同 workload/result fingerprint，减少顺序偏差。
核心测量区间覆盖查询、恢复写入、fsync 和复核；query block 单独报告纯查询区间。
这两个窗口不能混为一个“查询延迟”。

\code{make contest-demo} 自动构建 AgentOS Guest、启动 QEMU、解析串口并写出逐样本
raw、\code{measurements.csv}、\code{summary.json} 和报告。\code{dual-platform-run}
另构建 plain uCore 与 AgentOS-uCore，用确定性用户态合同做端到端比较。未实际运行时，
命令存在本身不构成性能数字。

\section{冻结 one-shot 主比较}

规范活动在同一源码提交上执行 16 个独立 workflow QEMU boot，平衡 AB/BA。16/16 个
core interval 中 indexed 更快；traversal/indexed core median 为 34.713/13.294 ms，
配对 speedup median 为 3.118 倍。

\ReportFigure[0.96\textwidth]
  {../figures/performance/07_core_end_to_end_scope.pdf}
  {../figures/performance/07_core_end_to_end_scope.pdf}
  {Core-path 与端到端窗口的配对结果及顺序分层}
  {fig:core-e2e-scope}

这组结果不能简化成“系统端到端快 3.118 倍”。端到端只有 3/16 个配对支持 indexed，
indexed minus traversal 的端到端差值中位数为 +13.452 ms，即当前 end-to-end 窗口中
indexed 更慢。AB 与 BA 两个顺序分层的 core speedup median 分别为 1.4535 倍与
5.4805 倍，说明顺序效应显著。我们因此只声称 core path latency 在该负载中改善，
并把端到端恢复、fsync、复核和测量噪声作为独立工程问题。

\ReportFigure[0.92\textwidth]
  {../figures/performance/01_paired_core_performance.pdf}
  {../figures/performance/01_paired_core_performance.pdf}
  {16 个独立 QEMU boot 的 traversal/indexed core 配对原始分布}
  {fig:paired-core}

\section{工作量与 I/O 口径}

文件查询图中的 \code{bytes_read} 是 workflow actor core-window lane work；其他 contest
numerator 是 shared global/end-to-end delta，再按 actor core-window observer syscall
归一。Task 路径则按 completed operation 归一。这些 owner、scope 和 window 不同的
数字不能混成一个进程归因，也不能跨 metric 用热力图颜色比较绝对量。

主文优先呈现配对 latency、catalog 网格、Task transport 与 EEVDF。混合 scope 的 I/O
热力图只放附录，并同时给 raw numerator、denominator、owner 和 window，避免漂亮的
颜色掩盖测量主体差异。

\section{与 Nexus 的职责分离}

Nexus 不替代 \code{contest-demo}。模型可能因用户问题、历史、失败或审批改变任务图，
这种动态性正是产品价值，却不适合作为等量性能基准。Nexus 可以读取历史
\code{nexus_meas} capsule 并在报告中引用已发布数字，但必须标记
\code{historical_not_this_boot}；它不会在启动时重跑 16 个 boot，也不能把 snapshot
改写成本次测量。

\section{验证}

Host runner 同时检查超时、panic、退出码、路径 fingerprint、AB/BA 样本完整性和输出
schema。one-shot manifest 记录源码 commit、Guest binary hash、工具版本、每个 boot
状态、33 个 raw 文件和 145-entry inventory；validation 要求 19 张长表 schema、数量、
来源 SHA 和 10 组 PNG/PDF 全部一致。规范目录写入 \code{COMPLETED} 后采集入口 fail
closed；离线重绘只读 CSV，不启动 QEMU。

\begin{evidencebox}[任务六结论]
我们交付了真实 QEMU 的多模块确定性综合程序和可复核性能数据，同时保留独立 Nexus
交互产品。当前最强性能结论是 indexed core path 在 16 个配对中全部更快；端到端结果
并不支持同等级加速主张。
\end{evidencebox}
